\documentclass[11pt]{article}
\usepackage[utf8]{inputenc}
\usepackage[T1]{fontenc}
\usepackage{lmodern}
\usepackage[margin=1in]{geometry}
\usepackage{amsmath,amssymb,mathtools}
\usepackage{graphicx}
\usepackage{grffile}
\usepackage{float}
\usepackage{hyperref}
\usepackage{caption}
\usepackage{parskip}
\usepackage{enumitem}
\usepackage{xurl}
\setlength{\parskip}{0.65em}
\setlength{\parindent}{0pt}
\begin{document}
\title{Daily Research Digest --- 2026-04-22}
\author{}
\date{}
\maketitle
\textbf{Topics:} galaxy clusters, galaxies, gravitational lensing, dark matter.
\section*{Synthesis}
Here is a summary of the key research topics covered by these papers:
1. **Dark Matter:**
\begin{itemize}
\item **"Extracting Dark-Matter Mass from Angular Scanning":** Proposes a novel method to determine dark matter mass using directionality observables in two-dimensional direct detection experiments.
\item **"Cosmological constraints on TeV-scale dark matter subcomponents decaying between recombination and reionisation":** Studies the impact of decaying TeV-scale dark matter particles on cosmological observables, with an emphasis on future radio telescopes' sensitivity to 21-cm signals.
\end{itemize}
2. **Galaxy-Galaxy Lensing:**
\begin{itemize}
\item **"Fourth-order galaxy-galaxy-lensing: Theoretical framework and direct estimation":** Extends the analytical basis for galaxy-galaxy lensing to fourth-order statistics, focusing on non-Gaussian features in the matter distribution. Develops methods for detecting these statistics from observations of Stage IV surveys.
\end{itemize}
3. **Gravitational Lensing:**
\begin{itemize}
\item **"SDSSJ110546.07+145202.4: The first long-duration radio changing-look NLS1 galaxy":** Describes a unique radio-changing look Narrow-line Seyfert 1 (NLS1) galaxy with high-frequency radio observations and detailed spectral energy distribution analysis.
\item **"Other red dots: A possible GLIMPSE of normal AGB stars at Cosmic Noon through extreme lensing":** Reports the discovery of extremely faint red point sources in JWST/NIRCam images, suggesting that these may be magnified asymptotic giant branch (AGB) stars or yellow super-giants.
\end{itemize}
4. **Galaxy Formation and Evolution:**
\begin{itemize}
\item **"Characterizing and spectrally modeling embedded FUor eruptions in the near-infrared":** Uses infrared surveys to identify and characterize FU Ori outbursts, providing detailed disk modeling with high-resolution spectroscopy.
\item **"SDSSJ110546.07+145202.4: The first long-duration radio changing-look NLS1 galaxy":** Also touches on the nature of AGN and accretion rate changes, contributing to our understanding of active galactic nuclei.
\end{itemize}
5. **Interstellar Medium (ISM) and High-Energy Phenomena:**
\begin{itemize}
\item **"Cosmological constraints on TeV-scale dark matter subcomponents decaying between recombination and reionisation":** Considers the impact of decaying dark matter particles on the ISM, particularly in the context of cosmic microwave background (CMB) observations.
\end{itemize}
6. **Astrophysical Phenomena:**
\begin{itemize}
\item **"Fourth-order galaxy-galaxy-lensing: Theoretical framework and direct estimation":** Investigates higher-order statistical measures to probe non-Gaussian features in the matter distribution.
\item **"Other red dots: A possible GLIMPSE of normal AGB stars at Cosmic Noon through extreme lensing":** Discusses how extreme gravitational lensing can reveal stellar populations that are otherwise too faint to observe.
\end{itemize}
These papers collectively contribute to our understanding of dark matter, galaxy formation and evolution, high-energy astrophysical phenomena, and the use of advanced observational techniques like gravitational lensing.
\newpage
\section*{Paper Summaries and Figures}
\subsection*{1. Precision Kinematic Sunyaev--Zel'dovich Measurements Across Halo Mass and Redshift with DESI DR2 and ACT DR6: Part II. Bright Galaxy Survey and Emission-Line Galaxies}
\textbf{Focus:} galaxies, gravitational lensing\\
\textbf{Authors:} B. Hadzhiyska, S. Ferraro, F. J. Qu, B. Ried Guachalla, E. Schaan, J. Aguilar, S. Ahlen, D. Bianchi, D. Brooks, F. J. Castander, E. Chaussidon, T. Claybaugh, A. de la Macorra, Arjun Dey, Biprateep Dey, P. Doel, J. E. Forero-Romero, E. Gaztañaga, S. Gontcho A Gontcho, G. Gutierrez, J. Guy, K. Honscheid, C. Howlett, D. Huterer, M. Ishak, R. Joyce, R. Kehoe, T. Kisner, A. Kremin, O. Lahav, M. Landriau, L. Le Guillou, A. Leauthaud, M. Manera, P. Martini, A. Meisner, R. Miquel, S. Nadathur, N. Palanque-Delabrouille, W. J. Percival, F. Prada, I. Pérez-Ràfols, G. Rossi, L. Samushia, E. Sanchez, E. F. Schlafly, D. Schlegel, J. Silber, D. Sprayberry, G. Tarlé, B. A. Weaver, R. Zhou, H. Zou\\
\textbf{Published:} 2026-04-21T17:59:30+00:00\\
\textbf{PDF:} \href{https://arxiv.org/pdf/2604.19745v1}{https://arxiv.org/pdf/2604.19745v1}\\
\textbf{Summary:}
We present the first high-significance spectroscopic stacked kinetic Sunyaev-Zel'dovich (kSZ) measurements of circumgalactic gas profiles for both Bright Galaxy Survey (BGS) and Emission Line Galaxy (ELG) tracers, combining DESI Data Release 2 with ACT Data Release 6. Using reconstructed line-of-sight velocities from the DESI galaxies and high-resolution ACT temperature maps, we detect the kSZ signal at high significance, reaching signal-to-noise ratios of up to $\sim$9 for BGS and $\sim$7.5 for ELGs in optimal stellar-mass selections. Together with the LRG measurements presented in Paper I, these constitute the most significant kSZ detections from any spectroscopic survey to date. We perform the analysis in both real and harmonic space, obtaining consistent results. By splitting both tracers into stellar-mass bins, we study the scaling of the kSZ amplitude with galaxy properties. Combining the kSZ measurements with ACT Data Release 6 (DR6) CMB lensing maps enables a joint calibration of the galaxy-halo connection and the gas fractions of host halos. For the BGS galaxies, we observe low gas fractions around the virial radius relative to standard expectations, likely attributable to active galactic nuclei (AGN) activity. We find some evidence for higher-mass halos retaining a larger fraction of their baryons, consistent with more efficient feedback in lower-mass systems. For the ELG sample, dominated by blue, star-forming galaxies, we provide the first detection of the gas distribution in ELG host halos. The ELGs appear to exhibit relatively high gas fractions, which points to the possibility of weaker feedback (due to e.g. low AGN and supernova feedback activity) at their mass scale. Finally, we present generalized Navarro-Frenk-White (GNFW) fits to the harmonic-space measurements, providing a compact parametrization of gas profiles for forward modeling in large-scale structure analyses.
\subsubsection*{Selected figures}
\begin{figure}[H]
\centering
\includegraphics[width=\linewidth,height=0.42\textheight,keepaspectratio]{/home/kf/research-assistant/data/cache/figures/http-arxiv.org-abs-2604.19745v1/figure\_1.png}

\end{figure}
\newpage
\subsection*{2. Precision Kinematic Sunyaev--Zel'dovich Measurements Across Halo Mass and Redshift with DESI DR2 and ACT DR6: Part I. Luminous Red Galaxies}
\textbf{Focus:} galaxies, dark matter\\
\textbf{Authors:} F. J. Qu, B. Ried Guachalla, E. Schaan, B. Hadzhiyska, S. Ferraro, J. Aguilar, S. Ahlen, A. Baleato Lizancos, D. Bianchi, D. Brooks, R. Canning, F. J. Castander, E. Chaussidon, T. Claybaugh, A. Cuceu, A. de la Macorra, B. Dey, P. Doel, A. Font-Ribera, J. E. Forero-Romero, E. Gaztañaga, S. Gontcho A Gontcho, G. Gutierrez, H. K. Herrera-Alcantar, K. Honscheid, C. Howlett, D. Huterer, M. Ishak, R. Kehoe, T. Kisner, A. Kremin, O. Lahav, M. Landriau, L. Le Guillou, M. E. Levi, M. Manera, A. Meisner, R. Miquel, S. Nadathur, J. A. Newman, W. J. Percival, I. P'erez-R`afols, G. Rossi, L. Samushia, E. Sanchez, E. F. Schlafly, D. Schlegel, M. Schubnell, H. Seo, J. Silber, D. Sprayberry, G. Tarl'e, B. A. Weaver, R. Zhou\\
\textbf{Published:} 2026-04-21T17:59:26+00:00\\
\textbf{PDF:} \href{https://arxiv.org/pdf/2604.19744v1}{https://arxiv.org/pdf/2604.19744v1}\\
\textbf{Summary:}
We present the most precise measurements of the kinetic Sunyaev-Zel'dovich (kSZ) effect around luminous red galaxies to date, detecting the signal at $18\sigma $ significance in both harmonic and configuration space. Our analysis cross-correlates 2.4 million spectroscopic LRGs from the Dark Energy Spectroscopic Instrument (DESI) DR2 sample with Data Release 6 (DR6) of the Atacama Cosmology Telescope (ACT). We develop a novel harmonic-space cross-correlation approach using momentum-weighted kSZ templates, yielding nearly uncorrelated bandpowers within a framework consistent with other large-scale structure analyses. By incorporating the LRG halo occupation distribution (HOD) and its uncertainty, we convert measured galaxy gas profiles into halo gas profiles and provide generalized Navarro-Frenk-White (GNFW) fitting profiles, providing empirical targets for tuning feedback efficiency in hydrodynamical simulations and for baryonic modeling in large-scale structure analyses. We find strong evidence that gas profiles do not trace dark matter, providing direct evidence for gas redistribution beyond gravitational collapse. Comparing to hydrodynamical simulations, our measurements favor feedback efficiencies exceeding those in the Battaglia profile, suggesting more efficient gas ejection in group-scale halos than previously predicted. Splitting by redshift, we detect the kSZ signal at SNR $\approx 5$--$10$ in each of four bins and find amplitude evolution consistent with the expected decline in mean halo mass at fixed comoving number density. Splitting by stellar mass, we study the scaling of kSZ amplitude with galaxy properties. Together with BGS and ELG measurements in Paper II, these results span $0.1 \lesssim z \lesssim 1.6$ across three galaxy populations, demonstrating the potential of spectroscopic kSZ to map circumgalactic gas and constrain baryonic feedback.
\subsubsection*{Selected figures}
\begin{figure}[H]
\centering
\includegraphics[width=\linewidth,height=0.42\textheight,keepaspectratio]{/home/kf/research-assistant/data/cache/figures/http-arxiv.org-abs-2604.19744v1/figure\_1.png}
\captionsetup{labelformat=empty}\caption{5.
Top: Projected distribution of DESI LRG galaxies in a representative sky region (120◦< RA < 160◦, 0◦<
. Each point represents an individual galaxy, with denser regions appearing darker due to overlapping markers
y purposes, the marker size (∼2 arcmin) is larger than the true angular extent of ind}
\end{figure}
\newpage
\subsection*{3. Breaking the UV Luminosity Function Degeneracy:Self-Interacting Dark Matter Constraints from Reionization Topology}
\textbf{Focus:} galaxies, dark matter\\
\textbf{Authors:} Zihan Wang, Huanyuan Shan\\
\textbf{Published:} 2026-04-21T17:50:29+00:00\\
\textbf{PDF:} \href{https://arxiv.org/pdf/2604.19726v1}{https://arxiv.org/pdf/2604.19726v1}\\
\textbf{Summary:}
Self-interacting dark matter (SIDM) is the leading framework resolving small-scale cold dark matter (CDM) crises, yet high-redshift SIDM constraints are fundamentally limited by degeneracies between dark matter microphysics and galaxy formation astrophysics. We demonstrate that the UV luminosity function alone cannot constrain SIDM: star formation suppression from SIDM halo core formation is fully absorbed by modest adjustments to standard astrophysical parameters. We show that 21 cm reionization topology breaks this degeneracy completely, providing a nuisance-immune probe: the SIDM-enhanced duty cycle of ionizing photon escape leaves a morphological signature fully independent of star formation efficiency. Combining JWST UVLF measurements with SKA1-Low forecasts, constant-cross-section SIDM with $\sigma /m \gtrsim 1$--$2\ \mathrm{cm^2/g}$ is either excluded or detectable across all physically motivated star formation coupling strengths. Our results establish a robust new avenue to probe dark matter microphysics in the early Universe.
\newpage
\subsection*{4. Resolved UV-Optical HST Imaging and Spectral Energy Distribution Modeling of Nearby BAT Active Galactic Nuclei}
\textbf{Focus:} galaxies\\
\textbf{Authors:} Connor Auge, Michael Koss, Kriti K. Gupta, Claudio Ricci, Benny Trakhtenbrot, Franz E. Bauer, Ezequiel Treister, Alessandro Peca, Brad Cenko, Kohei Ichikawa, Arghajit Janna, Darshan Kakkad, Richard Mushotzky, Kyuseok Oh, Alejandra Rojas Lilayú, David Sanders, Roberto Serafinelli, Matilde Signorini, Alessia Tortosa, C. Megan Urry\\
\textbf{Published:} 2026-04-21T16:54:19+00:00\\
\textbf{PDF:} \href{https://arxiv.org/pdf/2604.19674v1}{https://arxiv.org/pdf/2604.19674v1}\\
\textbf{Summary:}
We use high-resolution UV-to-optical imaging from the Hubble Space Telescope (HST) to construct spatially resolved spectral energy distributions (SEDs) for seven nearby ($z<0.07$) hard (14--195$\,$keV) X-ray-selected broad-line active galactic nuclei (AGN) with $L_{\rm bol}=10^{43.26}-10^{45.34}\,\rm{erg\,s^{-1}}$. The high spatial resolution of HST, which physically resolves structures on the scale of $\sim$50$\,$pc at $z=0.05$, enables the separation of AGN and host-galaxy emission through morphological decomposition with GALFIT, yielding improved measurements of AGN properties compared to those obtained with lower-resolution Swift UV/Optical Telescope (UVOT) data. AGN UV magnitudes derived from HST imaging (e.g., F225W) can differ by more than a magnitude from those from Swift/UVOT UVM2 due to extended nuclear emission. Additionally, the inclusion of high-resolution data at longer wavelengths (e.g., F814W) can significantly affect the resulting SED fit. Comparing fits of accretion disk and extinction models using HST and Swift/UVOT data, we find significant differences in the resulting parameters, with average differences of 2.0$\,$eV in the maximum disk temperature and 2.2$\,$mag in the AGN host-galaxy extinction. These differences ultimately lead to significant changes in bolometric luminosities and X-ray bolometric corrections, with the HST-based fits yielding average increases of $\sim$0.57$\,$dex and $\sim$0.66$\,$dex respectively. This demonstrates host-galaxy contamination in unresolved UV--optical data can strongly bias SED-based estimates of disk temperatures, extinction, bolometric luminosities, and X-ray bolometric corrections in AGN. Large-area, high-resolution imaging surveys from Euclid and the Nancy Grace Roman Space Telescope will extend these techniques to much larger AGN samples, enabling uniform, high-precision SED measurements in the near-IR.
\subsubsection*{Selected figures}
\begin{figure}[H]
\centering
\includegraphics[width=\linewidth,height=0.42\textheight,keepaspectratio]{/home/kf/research-assistant/data/cache/figures/http-arxiv.org-abs-2604.19674v1/figure\_1.png}
\captionsetup{labelformat=empty}\caption{ALFIT analysis used to constrain the photometric errors of the best fit point source
478 (left: F225W middle: F435W right: F814W. This shows the point source magnitude C. UV LIGHT CURVES e that variability may play in driving the differences in the measure point source
T data and HST data, we constr}
\end{figure}
\newpage
\subsection*{5. A Possible Protocluster of Galaxies Serendipitously Discovered in the Field of an Intermediate-Redshift Post-starburst Galaxy}
\textbf{Focus:} galaxies\\
\textbf{Authors:} Mary C. Knowlton, Justin S. Spilker, Rachel Bezanson, Vincenzo R. D'Onofrio, Anika Kumar, David J. Setton, Katherine A. Suess\\
\textbf{Published:} 2026-04-21T16:41:05+00:00\\
\textbf{PDF:} \href{https://arxiv.org/pdf/2604.19651v1}{https://arxiv.org/pdf/2604.19651v1}\\
\textbf{Summary:}
We present the serendipitous discovery of an overdensity of submillimeter galaxies (SMGs) in the field of SDSSJ0909-0108, a massive z\textasciitilde{}0.7 post-starburst galaxy from the SQuIGGLE survey. ALMA observations at 870um and 2mm reveal six galaxies within a 35'' region with flux ratios consistent with emission from dust. Given the rarity of 870um sources and the small field-of-view of ALMA, we speculate that some of these sources are physically associated. None of the sources are at the same redshift as the post-starburst, and four do not have spectroscopic redshifts. We suggest that follow-up optical and/or ALMA observations be carried out to measure redshifts for the galaxies in this potential protocluster environment.
\subsubsection*{Selected figures}
\begin{figure}[H]
\centering
\includegraphics[width=\linewidth,height=0.42\textheight,keepaspectratio]{/home/kf/research-assistant/data/cache/figures/http-arxiv.org-abs-2604.19651v1/figure\_1.png}
\captionsetup{labelformat=empty}\caption{Figure 1.
Top left: ALMA Band 4 continuum image.
The white box indicates the FOV of the Band 7 image.
All six
sources are labeled with their primary-beam-corrected fluxes in parentheses.
Top right: ALMA Band 7 continuum image,
with primary-beam-corrected flux densities in parentheses for the five vi}
\end{figure}
\newpage
\subsection*{6. MG-NECOLA: A Field-Level Emulator for $f(R)$ Gravity and Massive Neutrino Cosmologies}
\textbf{Focus:} galaxies\\
\textbf{Authors:} J. Bayron Orjuela-Quintana, Mauricio Reyes, Elena Giusarma, Marco Baldi, Neerav Kaushal, César A. Valenzuela-Toledo\\
\textbf{Published:} 2026-04-21T16:01:37+00:00\\
\textbf{PDF:} \href{https://arxiv.org/pdf/2604.19613v1}{https://arxiv.org/pdf/2604.19613v1}\\
\textbf{Summary:}
Accurate modeling of non-linear gravitational dynamics is essential for constraining extensions to the standard cosmological model using large-scale structure observations. While high-resolution $N$-body simulations provide the required fidelity, they are computationally prohibitive for the large ensembles needed to analyze Modified Gravity (MG) scenarios. We present MG-NECOLA, a field-level emulator based on a convolutional neural network that upgrades fast, approximate MG-PICOLA simulations to near--$N$-body accuracy at a fraction of the computational cost. Trained on a suite of QUIJOTE\_MG simulations for $f(R)$ gravity, MG-NECOLA achieves nearly sub-percent accuracy ($\lesssim 1\%$) in both the matter power spectrum and bispectrum up to $k \simeq 1~h\,\mathrm{Mpc}^{-1}$. Crucially, although being trained on a fixed cosmology, the network generalizes robustly to cosmologies outside its training manifold keeping the error below $5\%$. It successfully recovers the General Relativity limit ($Λ$CDM) without introducing spurious MG signals and accurately captures the power suppression induced by massive neutrinos ($M_ν\leq 0.4$ eV), despite being trained on cosmologies with massless neutrinos. The pipeline delivers a speed-up factor of $\sim 1500\times$ relative to full $N$-body runs, generating a high-fidelity realization in O$(10^3)$ CPU seconds compared to O$(10^6)$ for the baseline. This accuracy-efficiency trade-off establishes MG-NECOLA as a powerful tool for generating the massive mock catalogs required for next-generation galaxy surveys.
\subsubsection*{Selected figures}
\begin{figure}[H]
\centering
\includegraphics[width=\linewidth,height=0.42\textheight,keepaspectratio]{/home/kf/research-assistant/data/cache/figures/http-arxiv.org-abs-2604.19613v1/figure\_1.png}
\captionsetup{labelformat=empty}\caption{Figure 1. Schematic representation of the V-Net architecture adopted in this work. The network follows an encoder--decoder
(contracting--expansive) design with two down-sampling and two up-sampling stages, connected through residual and skip
connections at multiple resolutions. The input to the networ}
\end{figure}
\newpage
\subsection*{7. Understanding supernova gravitational waves with protoneutron star asteroseismology}
\textbf{Focus:} galaxies\\
\textbf{Authors:} Hajime Sotani\\
\textbf{Published:} 2026-04-21T15:10:34+00:00\\
\textbf{PDF:} \href{https://arxiv.org/pdf/2604.19557v1}{https://arxiv.org/pdf/2604.19557v1}\\
\textbf{Summary:}
Supernovae are one of the most promising gravitational wave sources. But, since the system of the supernovae is nearly spherically symmetric, the expected gravitational waves from them are relatively weak, compared to the case of the compact binary mergers. Thus, at least using the current gravitational wave detectors, only the gravitational waves from a supernova that occurred in our galaxy could be detected. To reliably extract information from gravitational waves originating from such a low event rate, thorough preparation is essential. However, because supernova gravitational waves strongly depend on model parameters, such as progenitor mass and the equation of state for dense matter, it may be difficult to extract physical properties even if the gravitational waves are detected. The universal relations between gravitational-wave signals and physical properties, independent of model parameters, are important for solving this difficulty. To discuss such a universal relation, in this article, we systematically examine the protoneutron-star oscillation frequencies with the linear analysis, the so-called asteroseismology, and compare them with the gravitational wave signals in the simulations.
\newpage
\subsection*{8. Characterizing and spectrally modeling embedded FUor eruptions in the near-infrared}
\textbf{Focus:} galaxies\\
\textbf{Authors:} Jiaxun Li, Tinggui Wang, Zheyu Lin\\
\textbf{Published:} 2026-04-21T15:07:07+00:00\\
\textbf{PDF:} \href{https://arxiv.org/pdf/2604.19551v1}{https://arxiv.org/pdf/2604.19551v1}\\
\textbf{Summary:}
Context. Episodic accretion in young stellar objects (YSOs) is thought to play a critical role in addressing the "luminosity problem" associated with star formation. However, optical surveys tend to bias against sources that are heavily obscured. Infrared time-domain surveys, such as unTimely WISE, facilitate the identification of such sources within the dense star formation regions of our Galaxy. Aims. We aim to systematically identify and characterize FUor outbursts in infrared-selected YSOs using high-resolution spectroscopy and detailed disk modeling. Methods. We conducted follow-up high-resolution spectroscopy with Gemini South/IGRINS for four FUor candidates discovered in infrared time-domain surveys. Using a combination of photometric and spectroscopic observations, we constructed spectral energy distributions and fit them with a disk model that incorporates an actively accreting inner disk together with a passively irradiated outer disk. Results. All objects show CO and H$_2$O absorption bands at 2.3$\mu $m, and their positions in the Na + Ca versus CO equivalent width diagram further corroborate their classification as FUors. The best-fitting model spectra closely match both the observed spectral features and the overall continuum, providing additional confirmation of the FUor classification. The best-fit models reveal high extinction values ($A_V$ = 10-20 mag), with $M_*\dot{M}$ comparable to those of classical FUors such as FU Orionis. Among 18 sources initially selected via infrared light curves, $6-$7 out of 8 with available spectra exhibit FUor characteristics, implying a high selection efficiency.
\subsubsection*{Selected figures}
\begin{figure}[H]
\centering
\includegraphics[width=\linewidth,height=0.42\textheight,keepaspectratio]{/home/kf/research-assistant/data/cache/figures/http-arxiv.org-abs-2604.19551v1/figure\_1.png}
\captionsetup{labelformat=empty}\caption{Fig. B.1: Distribution of parameters and corner diagram for
V191. The upper panel shows the near solution, and the lower
panel shows the far solution. The markings in the figure are con-
sistent with Fig. 4.}
\end{figure}
\newpage
\subsection*{9. SDSSJ110546.07+145202.4: The first long-duration radio changing-look NLS1 galaxy}
\textbf{Focus:} gravitational lensing\\
\textbf{Authors:} S. Komossa, D. Grupe, A. Kraus, P. G. Edwards, E. F. Kerrison, K. Rose, R. Soria, T. An, M. J. Hardcastle, K. E. Gabanyi, S. Panda, D. W. Xu, J. Wang, S. Frey, A. Mezosi\\
\textbf{Published:} 2026-04-21T13:09:40+00:00\\
\textbf{PDF:} \href{https://arxiv.org/pdf/2604.19435v1}{https://arxiv.org/pdf/2604.19435v1}\\
\textbf{Summary:}
SDSSJ110546.07+145202.4 stands out as a unique radio changing-look Narrow-line Seyfert 1 (NLS1) galaxy that has brightened dramatically and shows an exceptionally long duration of its "on" phase. We present the first high-frequency radio observations, the first simultaneous radio spectral energy distributions (SEDs), the first optical--UV--X-ray SEDs, and the first X-ray monitoring and spectroscopy of this recently discovered event. Importantly for understanding the nature of the outburst, we show that the X-ray spectrum is soft with a photon index Gamma\_X=2.5; line-of-sight absorption and extinction are low or absent; the radio SED is peaked at low frequencies \textasciitilde{}2 GHz; and the radio outburst emission is very long-lived (t > 8 yr) and roughly constant. The softness of the X-ray spectrum, low supermassive black hole (SMBH) mass, and high Eddington ratio all corroborate the optical NLS1 classification. We discuss multiple outburst scenarios, including lensing, absorption, a binary SMBH merger, a long-duration giant-star tidal disruption, a newly ignited active galactic nucleus (AGN), and an accretion-rate change. While most of them can be either excluded or are deemed too rare and lack positive evidence so far, most or all types of these transients are expected to be detected in ongoing VLA and upcoming SKA surveys. SDSSJ110546.07+145202.4 itself is well explained by an accretion rate change that triggered the powerful radio jet emission. The low redshift and SMBH mass of this system offer a unique perspective of the physical processes of radio-jet ignition that are expected to operate in the early Universe around growing SMBHs.
\subsubsection*{Selected figures}
\begin{figure}[H]
\centering
\includegraphics[width=\linewidth,height=0.42\textheight,keepaspectratio]{/home/kf/research-assistant/data/cache/figures/http-arxiv.org-abs-2604.19435v1/figure\_1.png}
\captionsetup{labelformat=empty}\caption{6. SMBH MASS ESTIMATES, EDDINGTON RATIO
AND NLS1 CLASSIFICATION}
\end{figure}
\newpage
\subsection*{10. Cosmological constraints on TeV-scale dark matter subcomponents decaying between recombination and reionisation}
\textbf{Focus:} dark matter\\
\textbf{Authors:} Markus R. Mosbech, Cristina Benso, Felix Kahlhoefer\\
\textbf{Published:} 2026-04-21T08:07:34+00:00\\
\textbf{PDF:} \href{https://arxiv.org/pdf/2604.19198v1}{https://arxiv.org/pdf/2604.19198v1}\\
\textbf{Summary:}
The Dark Ages and the Cosmic Dawn are an untapped well of information about the particle physics properties of dark matter, which may become accessible with future radio telescopes able to probe the 21-cm signal from atomic hydrogen. In this work we study the impact on cosmological observables of a dark matter subcomponent composed of TeV-scale particles that decay into electrons, photons or neutrinos with a lifetime shorter than the age of the universe. We re-evaluate constraints from the Cosmic Microwave Background (CMB) on these scenarios using the most recent data sets and estimate the sensitivity of future detections of the global 21-cm signal. Our main result is that the latter is potentially more sensitive to the effects of decaying dark matter with a lifetime $\tau \gtrsim 10^{15} \, \mathrm{s}$. This effect is strongest for the case of decays into neutrinos due to the different spectral distribution of the injected electromagnetic energy. For DM masses well above the TeV-scale, these differences become less pronounced and the sensitivity of both the CMB and the 21-cm signal depend primarily on the total amount of injected electromagnetic energy.
\newpage
\subsection*{11. Extracting Dark-Matter Mass from Angular Scanning}
\textbf{Focus:} dark matter\\
\textbf{Authors:} Daeyeong Jeong, Doojin Kim, Jong-Chul Park\\
\textbf{Published:} 2026-04-20T18:03:24+00:00\\
\textbf{PDF:} \href{https://arxiv.org/pdf/2604.18704v1}{https://arxiv.org/pdf/2604.18704v1}\\
\textbf{Summary:}
We propose a novel method to determine the mass scale of ambient dark matter, applicable to (at least effectively) two-dimensional direct detection experiments that allow for directionality observables. Due to the motion of the solar system and Earth relative to the Galactic Center and the Sun, the dark-matter flux exhibits a directional preference. We first demonstrate that dark-matter event rates depend non-trivially on the angle between the detection plane and the overall dark-matter flow, with the curvature of this angular spectrum encoding mass information. As proof of principle, we take the recently proposed Graphene-Josephson-Junction-based superlight dark-matter detector as a concrete example and validate these theoretical expectations through numerical analyses.
\subsubsection*{Selected figures}
\begin{figure}[H]
\centering
\includegraphics[width=\linewidth,height=0.42\textheight,keepaspectratio]{/home/kf/research-assistant/data/cache/figures/http-arxiv.org-abs-2604.18704v1/figure\_1.png}
\captionsetup{labelformat=empty}\caption{FIG. 4. Flow chart of DarkWind.}
\end{figure}
\newpage
\subsection*{12. Other red dots: A possible GLIMPSE of normal AGB stars at Cosmic Noon through extreme lensing}
\textbf{Focus:} gravitational lensing\\
\textbf{Authors:} Lukas J. Furtak, Adi Zitrin, Erik Zackrisson, Vasily Kokorev, Anthony J. Taylor, Joseph F. V. Allingham, John Chisholm, Jose M. Diego, Hakim Atek, Kristen B. W. McQuinn, Ryan Endsley, Richard Pan, Gabriel Brammer, Qinyue Fei, Seiji Fujimoto, Tiger Y. -Y. Hsiao, Patrick L. Kelly, Damien Korber, Ashish K. Meena, Rohan P. Naidu, Alberto Saldana-Lopez\\
\textbf{Published:} 2026-04-20T18:00:10+00:00\\
\textbf{PDF:} \href{https://arxiv.org/pdf/2604.18696v1}{https://arxiv.org/pdf/2604.18696v1}\\
\textbf{Summary:}
We report the discovery of four extremely faint ($m_{\mathrm{F444W}}\gtrsim29$) red point sources in recent ultra-deep JWST/NIRCam images of the strong lensing galaxy cluster Abell S1063. All four sources sit in lensed arcs, on the symmetry points very close to the critical curves for their host-galaxies' redshifts ($z\sim1-4$). Remarkably, these point sources appear in most arcs that are sufficiently faint close to the critical curve's position ($<21\,\mathrm{nJy}\,\mathrm{arcsec}^{-2}$ in F115W). This suggests that -- unlike previous caustic-crossing events or lensed stars -- thanks to the unprecedented depth of the GLIMPSE observations paired with the extreme lensing magnification (up to $\mu \sim10^4$) we might be resolving the lower-mass ($M\sim1-11\,\mathrm{M}_{\odot}$) red stellar population. Concretely, we detect three likely extremely magnified asymptotic giant branch (AGB) stars ($T_{\mathrm{eff}}\sim3200-3750$ K), and one yellow super-giant star ($T_{\mathrm{eff}}\sim6750$ K) -- possibly a yellow hyper-giant or a Cepheid. In addition to offering the first glimpse at low-mass extremely magnified stars, these detections open a possible window into stellar populations, evolution, and chemical enrichment at high redshifts, and could pave the way for using lensed stars such as these as standard candles to populate the distance ladder at cosmological redshifts.
\subsubsection*{Selected figures}
\begin{figure}[H]
\centering
\includegraphics[width=\linewidth,height=0.42\textheight,keepaspectratio]{/home/kf/research-assistant/data/cache/figures/http-arxiv.org-abs-2604.18696v1/figure\_1.png}
\captionsetup{labelformat=empty}\caption{Figure 2. Simulated brightness properties of a 4.5 Gyr old, constant SFH stellar population of M⋆= 6 × 105 M⊙at z = 1.43
(the spectroscopic redshift of ID 24951). This represents the most conservative population in the vicinity of the lensed star
since it corresponds to the age of the Universe at th}
\end{figure}
\newpage
\subsection*{13. Fourth-order galaxy-galaxy-lensing: Theoretical framework and direct estimation}
\textbf{Focus:} gravitational lensing\\
\textbf{Authors:} Jonathan Oel, Lucas Porth, Peter Schneider, Elena Silvestre-Rosello\\
\textbf{Published:} 2026-04-20T15:04:01+00:00\\
\textbf{PDF:} \href{https://arxiv.org/pdf/2604.18378v1}{https://arxiv.org/pdf/2604.18378v1}\\
\textbf{Summary:}
Traditional galaxy-galaxy lensing is a well-established method of probing the statistical properties of the Universe's matter and galaxy distribution. However, this measure does not carry all the statistical information, provided the matter and galaxy distribution contain non-Gaussian features. In order to study these non-Gaussianities, it is necessary to consider higher-order statistical measures. The aim of this work is to extend the analytical basis describing the statistical correlations between galaxies and shear to the fourth order, with special emphasis on the associated aperture statistics. In order to include fourth-order statistics in future analysis of the relation between mass and galaxies, we further investigate whether we can expect to detect these statistics from observations of stage IV surveys. We define the four-point correlation function (4PCF) between the shear and the positions of triplets of foreground galaxies and derive its relation to the respective trispectrum. We convert the 4PCF to aperture statistics and derive the analytical form of the respective filter function, which we then implement in a numerical integration pipeline. Furthermore, we develop a direct estimator that allows us to measure galaxy-mass aperture moments of arbitrary order on pixelized data using a Fast-Fourier-Transform (FFT) algorithm. We show that the corresponding aperture measure $\langle\mathcal{N}^3 M_\mathrm{ap}\rangle$ can be calculated with sub-percent accuracy on relevant aperture scales, $θ$, by means of numerical integration. Furthermore, we apply the FFT-based direct estimator to a mock catalog with a realistic stage IV survey setup on a sky area of $2000~\mathrm{deg}^2$, and detect the connected part of the aperture statistics $\langle\mathcal{N}^3 M_\mathrm{ap}\rangle(θ)$ with a signal-to-noise ratio of roughly nine on small aperture scales.
\subsubsection*{Selected figures}
\begin{figure}[H]
\centering
\includegraphics[width=\linewidth,height=0.42\textheight,keepaspectratio]{/home/kf/research-assistant/data/cache/figures/http-arxiv.org-abs-2604.18378v1/figure\_1.png}
\captionsetup{labelformat=empty}\caption{Fig. A.2. The filter function (ϑ1ϑ2ϑ3)2AN3M(ϑ1, ϑ2, ϑ3, ϕ12, ϕ23|θ) as a function of ϑ1,2/θ on a logarithmic scale for four different configurations of
(ϕ12, ϕ23). The top plots show the real part and the bottom plots the imaginary part of the function. All aperture scales are fixed to θ = 5′, as we}
\end{figure}
\end{document}
